\chapter{State of the Art}
L'evoluzione tecnologica ha permesso un nuovo paradigma per affrontare le sfide che il cambiamento climatico pone a tutti noi, questo nuovo modo di affrontare i problemi sul cambiamento climatico si basa su un paradigma data-driven; ciò è possibile grazie alla disponibilità di enormi dataset storici che oggigiorno abbiamo e all'impressionante miglioramento delle capacità computazionali, il machine learning  a questo punto non è più solo un supporto, grazie a tutto ciò è diventato uno strumento essenziale nelle sfide che il riscaldamento globale pone di fronte a tutti noi.\\
Le tecniche di machine learning sono diventate essenziali per l'adattamento ai cambiamenti climatici e per lo sviluppo sostenibile, in quanto esse forniscono capacità predittive fondamentali per permettere una pianificazione migliore, rendendoci in grado di perseguire uno sviluppo più sostenibile, inoltre esse sono largamente adattabili alle varie zone del mondo nelle quali vengono applicate, ed infine, Grazie alla sempre maggiore disponibilità di nuovi dati, è possibile progredire apprendendo continuamente da essi. Questo approccio è di particolare utilità nel contesto del cambiamento climatico, dove le condizioni si evolvono continuamente nel tempo; permette quindi ai modelli di machine learning di migliorare la loro accuratezza man mano che incontrano nuovi dati, garantendo strategie di adattamento aggiornate ed efficaci.

\section{Ambienti d'uso del ML}
Come detto le tecniche di machine learning possono essere molto utili nell'interpretare enormi quantità di dati ambientali, che oggigiorno vengono prodotti ad una velocità e con una frequenza maggiore che in passato (Faghmous Kumar et al.~\cite{faghmous2014bigdata}, 2014), in modo tale da vedere relazioni non lineari e che sarebbero difficilmente visibili senza l'utilizzo del machine learning, permettendo ai sistemi climatici di diventare più efficaci e rendendoci capaci di creare strategie atte all'ottimizzazione delle risorse. Grazie alla grande versatilità dei sistemi di apprendimento automatico, essi possono essere applicati nei campi più disparati, bisogna però tenere a mente che le tecniche di machine learning possono sì far quanto detto aumentando la capacità di monitoraggio di vari parametri, velocizzando ed automatizzando il processo di elaborazione dei dati raccolti, possono inoltre in generale velocizzare il processo scientifico e l'efficienza delle varie soluzioni possibili, però essi rimangono solo una parte della soluzione che può e deve essere integrato ad altri strumenti e processi per raggiungere dei risultati concreti (Rolnick et al. ~\cite{rolnick2022tackling}, 2022)

\subsection{Applicazioni per la previsione nei cambiamenti climatici}
Un tema fondamentale nell'applicazione di strumenti di machine learning nella lotta e adattabilità delle attività umane al cambiamento climatico è la capacità di questa tecnologia di effettuare delle previsioni climatiche accurate e tempestive, ciò è incredibilmente utile perchè attraverso queste tecniche si riesce ad elaborare complessi e grandi quantità di dati, come per esempio informazioni ricevute da immagini satellitari e dai sensori climatici continuamente, riuscendo a rilevare correlazioni che non sono immediatamente riconoscibili, cosa molto comune nelle rappresentazioni del sistema terrestre essendo estremamente complesso(Huntingford,~\cite{huntingford2019machine} 2019), ciò ci permette, grazie a queste previsioni, di agire meglio e più efficacemente che senza il loro utilizzo, rendendoci in grado di identificare i segnali di un'ipotetica degradazione ambientale in anticipo, cosi da poter agire tempestivamente preparandoci agli scenari climatici futuri che ci attendono.\\
In particolare per esempio i modelli di deep learning, hanno dimostrato come essi siano una risorsa fondamentale nel catturare le complesse relazioni tra i dati climatici, grazie alla loro vasta adattabilità e comprovata accuratezza nei risultati.\\

\subsection{Appllicazioni riguardanti la sostenibilità ambientale}
Gli algoritmi di machine learning possono essere utilizzare per raggiungere una maggiore sostenibilità ambientale garantendo uno sviluppo efficiente e non distruttivo nei confronti dell'ambiente intorno a noi, migliorando l'ottimizzazione delle risorse, riducendo lo spreco e l'impatto ambientale, ma più in generale per esempio fornendo agli agricoltori più informazioni, permettendo un minore spreco di risorse quali acqua, fertilizzanti, pesticidi etc, in generale i modelli di machine learning permettono "l'evoluzione dei sistemi di gestione agricola in veri e propri sistemi di intelligenza artificiale" (Liakos,~\cite{liakos2018machine} 2018).\\
Oltre che per garantire uno sviluppo più sostenibile gli strumenti basati su tecniche di machine learning possono anche essere usati per aiutare a monitorare l'ambiente con la sua fauna e flora, per esempio questi tecniche possono essere utilizzate insieme ad algoritmi di riconoscimento delle immagini per identificare specie diverse da foto trappole, riducendo il bisogno di controllo manuale delle immagini (Norouzzadeh et al. ~\cite{norouzzadeh2018automatically} 2018), in modo simile i modelli di machine learning possono essere usati nel contesto dell'economia circolare per riconoscere materiali riciclabili e non riciclabili, o ancora per rivelare tempestivamente cambi nei pattern climatici e delle precipitazioni, applicandole all'output di grandi modelli climatici.\\

\section{Casi presi in esame nell'adattamento ai cambiamenti climatici attraverso tecniche di machine learning}

Le tecniche di machine learning sono state applicate in vario modo nei campi citati e anche in altri, per mostrare ciò di seguito si analizzeranno vari studi che mostreranno alcune delle apllicazioni reali possibili che sono state ipotizzate:

\subsection{Previsione del deficit di pressione di vapore per la gestione delle risorse idriche agricole mediante apprendimento automatico in ambienti semi-aridi}
Nel caso preso in esame da Elbeltagi et al.~\cite{elbeltagi2023vpd} vengono applicati modelli di machine learning per permettere la previsione del deficit di pressione di vapore (VPD), un parametro chiave che influenza il calcolo dell’evatraspirazione(ET), valore fondamentale per la gestione delle risorse idriche in agricoltura.
Lo studio si concentra su 8 regioni del delta del nilo (Dakahliyah, Gharbiyah, Kafr Elsheikh, Dumyat, Port Said, Ismailia, Sharqiyah e Qalubiyah).
Per perseguire quest'obiettivo sono stati impiegati sei algoritmi di apprendimento automatico(ML): linear regression(LR), additive regression tree (ART), random subspace(RSS), random forest (RF), algoritmo M5 di Quinlan (M5P), reduced error pruning tree (REPTree), in questo studio è risultato come modello migliore il random forest, mostrando delle ottime prestazioni sia nella fase di test che in quella di training, ovvero un coefficiente di correlazione di 0,9694,un errore assoluto medio di 0,0967, un errore quadratico medio di 0,1252, percentuali di errore relativo del 21,7297 e di errore quadratico medio del 24,0356, dimostrando una grande robustezza nelle previsioni del VPD. Questi risultati mostrano come il modello random forest si sia dimostrato uno strumento molto utile per gli studi idro-climatologici.

\subsection{Predictive modeling of wildfires: A new dataset and machine learning approach}
L'articolo di Younes Oulad Saya et al.~\cite{ouladsaya2023wildfires} si concentra sul dare una risposta a queste tre domande:
\begin{itemize}
	\item Come possiamo costruire un dataset basato su dati di telerilevamento?
	\item Come possiamo processare il nostro dataset con algoritmi di data mining?
	\item Come possiamo migliorare la prevedibilità del nostro modello?
\end{itemize}
Il primo passo per costruire il dataset è stato quello di raccogliere dati che possano dimostrare la comparsa di un incendio, per fare ciò sono stati selezionati tre parametri, salute delle colture, temperatura del suolo e un indicatore di incendio, nello specifico è stato usato l'Indice di vegetazione a differenza normalizzata (NDVI), la temperatura della superficie terrestre ed infine le anomalie termiche.\\
Una volta che i dati sono stati collezionati sono state eseguite varie operazioni di prepocesso sui dati per verificare potenziali imperfezioni in essi.
\begin{itemize}
	\item Innazitutto i dati sono stati convertiti in formati leggibili.
	\item Una volta converti, sono stati "puliti" per correggere eventuali imperfezioni nei dati grezzi. 
	\item successivamente alla pulizia dei dati è stato effettuata una selezione delle immagini satellitari per includere unicamente le aree che hanno subito incendi.
	\item Generati i dati raster delle aree bruciate, si è passati alla normalizzazione dei dati mantenendo un unico timestamp, per fare ciò è stata eseguita un'interpolazione spazio-temporale dei dati.
	\item effetuata l'interpolazione spazio-temporale dei dati, si deve passare all'estrapolazione dei dati spazio-temporali, per fare ciò vengono realizzate delle stime sui dati ottenuti per prevedere incendi futuri.
	\item infine l'ultimo passaggio nella creazione del datset è stato quello dell'estrazione dei parametri utili dalle immagini satellitari.
\end{itemize}
Per predire l'avvento di incendi, si utilizzeranno algoritmi di data mining, nello specifico in quest'articolo sono stati usati reti neurali (MLP) e SVM.\\
L'area di studio monitorata nell'articolo è localizzata al centro del Canada (per la maggior parte British Columbia e Quebec), essa è stata scelta per l'elevato tasso di incendi boschivi e per analizzare i dati è stata usata la piattaforma per analisi di big data Databricks.\\
per effettuare lo studio sono state selezionate zone di incendi boschivi avvenuti tra il 2013 ed il 2014, per ognuna di queste zone sono stati usati i dati descritti in precedenza.\\
I risultati delle previsioni hanno mostrato come entrambi gli algoritmi hanno fornito buoni risultati con percentuali di accuratezza attestate al 97.48 per SVM e 98.32 per la rete neurale; Un confronto con alcuni modelli di previsione degli incendi boschivi ha inoltre dimostrato le prestazioni di questo modello. Tutti questi risultati confermano l'efficacia del modello nel prevedere il verificarsi di incendi boschivi.

\subsection{Machine Learning-Based Temperature Forecasting for Sustainable Climate Change Adaptation and Mitigation}

Nella fattispecie questo studio effettuato da Fatih Sevgin.~\cite{sevgin2023temperature}  vuole dimostrare come sia possibile sviluppare un modello per la previsione della temperatura ad Istanbul sfruttando reti neurali artificiali, comparando poi i risultati riscontrati da esse anche con altri modelli di machine learning, nello specifico vuole fornire un contributo nel mostrare le sfide che la creazione di previsioni meteorologiche a lungo termine comportino.\\
Nella prima parte dello studio è stata usata una rete neurale artificiale feedforward(ANN) per prevedere la temperature partendo da parametri rappresentanti l'umidità, la velocità del vento e le precipitazioni. Il modello ANN era costituito da 20 livelli nascosti con dati normalizzati attraverso la funzione di minmax, i risultati forniti dal modello hanno mostrato come le stime prodotte mostrassero dei risultati estremamente soddisfacenti.\\
Nella seconda parte dello studio sono stati utilizzati altri modelli di machine learning: modelli lineari (LM), support vector machine (SVM), K-nearest neighbor (KNN) e random forest (RF). Nell'utilizzare questi modelli lo studio ha utilizzato la tecnica della "cross-validation" per rendere le previsioni del dataset più accurate e realistiche, verificando i risultati attraverso il MAE, RMSE, and R-squared metric si vede innanzitutto come all'aumentare del numero di k-fold aumenta l'accuratezza del modello, inoltre comparando i risultati dei vari modelli si vede come il random forest si dimostri in generale il modello che fornisce le performance più bilanciate, raggiungendo un'accuratezza abbastanza alta sia con un numero elevato di k-fold, ma seppur peggiore, anche con un numero più contenuto degli stessi.\\
Questo studio mostra dunque come modelli di reti neurali artificiali e di machine learning siano una risorsa affidabile nell'affrontare le sfide che le previsioni meteo a lungo termine possono comportare, essi, come visto nell'articolo, forniscono infatti livelli di accuratezza davvero elevati sia nel caso delle reti neurali artificiali che nei vari modelli di machine learning provati, in particolare col Random Forest (RF)  mostrando come essi siano capaci di modellare la non linearità e dunque l'eterogeneità delle informazioni meteorologiche.

\subsection{Improving daily reference evapotranspiration forecasts: Designing AI-enabled recurrent neural networks based long short-term memory}

Lo studio di Ali et al. ~\cite{ali2023evapotranspiration} si pone come obiettivo quello di verificare quanto i nuovi modelli di deep learning riescano a comprendere le relazioni tra diversi tipi di dati in modo da riuscire ad effettuare delle previsioni accurate  sul valore dell'evapotraspirazione giornaliera (ETo), un parametro che ha un ruolo fondamentale nel ciclo idrologico in quanto da esso dipendono i cicli dell'acqua, dell'energia e del carbonio; E' fondamentale per esempio soprattutto in aree con scarse risorse idriche in quanto in particolare in queste zone influenza la totalità del bilancio idrico.\\
Per fare quanto citato lo studio prende in esempio i valori delle stazioni climatiche di due città australiane: Sydney e Melbourne.
\begin{itemize}
	\item Innanzitutto i dati sono stati preprocessati in modo da escludere i parametri poco o per nulla influenti sull'Eto, per verificare ciò si sono viste le correlazioni tra i vari parametri attraverso due heatmap, una per melbourne ed una per Sydney, fatto ciò i dati sono stati normalizzati per migliorare le performance dei modelli successivamente, i parametri scelti dopo ciò sono stati: la temperatura massima e minima giornaliera, l'umidità relativa,la radiazione solare giornaliera ed i lag dell'evapotraspirazione di riferimento
	\item Una volta scelte le feature il dataset è stato diviso in due parti, una costituita dall'80\% del totale come training set, ovvero i dati facenti riferimento al lasso di tempo 2009-2020, ed il restante 20\% per il test set, ovvero gli anni dal 2021 al 2024
	\item Non sono state effetuate tecniche di validazione come per esempio la cross-validation in quanto essendo i dati serie temporali non si voleva rompere l'ordine temporale degli stessi mettendo per esempio valori del 2021 nel training e valori del 2019 nel test-set
	\item infine per valutare i vari modelli testati sono state analizzate varie metriche: coefficiente di correlazione (R), coefficiente di determinazione (R2) ed errore quadratico medio (RMSE). L'analisi di queste metriche ha mostrato come l'algoritmo RNN-LSTM ha prodotto nettamente i risultati migliori in tutti e tre i criteri di valutazione in entrambe le stazione delle due città.
\end{itemize}
In conclusione lo studio ha mostrato come tutti i modelli testati hanno prodotto risultati soddisfacenti e come nello specifico il modello RNN-LSTM abbia avuto dei risultati incredibilmente positivi.

\subsection{Conclusione}
La varietà dei casi presi in esame attesta come le tecniche di machine learning siano uno strumento fondamentale nello studio e applicazione di soluzioni atte ad un possibile adattamento al cambiamento climatico grazie alle capacità esplicitate in precedenza quali l'adattamento alle varie zone climatiche del pianeta, come si è visto gli studi vanno dal freddo Nord America alla calda Australia, dimostrando in entrambi i casi ottimi risultati, e la capacità di elaborare enormi quantità di dati in continuo aumento, riuscendo per esempio senza nessun problema ad elaborare serie temporali di dati che partono addirittura dalla metà del secolo scorso fino ai giorni nostri, capacità che non è minimamente ai suoi limiti e che con l'aumento delle banche dati a disposizione renderà il lavoro degli algoritmi di machine learning sempre più essenziale e soprattutto utile, ed è in questo contesto che questa tesi si inserisce in quanto si ritiene che lo studio su algoritmi di machine learning non vada solo demonizzato nel contesto dello spreco di risorse ma vada dimostrato come si sta già facendo che esso possa essere anche utilizzato come uno strumento positivo nella lotta al cambiamento climatico permettendo analisi e trovando soluzioni che senza di esso erano troppo dispendiose o addirittura impossibili da apllicare

\label{sec:sota}